% =============================================================================
% Appendix: Supplementary Material
% =============================================================================
% SPDX-FileCopyrightText: 2026 RutherfordD@cardiff.ac.uk <frank@langbein.org>
% SPDX-License-Identifier: LPPL-1.3c
% =============================================================================

\section{File Structure}

FIG Shows the back-end file structure of CatCode. All addresses are relative to \texttt{res://src} (project source files). For clarity, when the \texttt{.gd} and\texttt{.tscn} for a node script pair are in the same folder, it will be listed as \texttt{.gd/tscn}. Not all files are listed here.

\begin{fullwidthfigure}[h]
    \centering
    \caption{Back-end file structure of the catCode systems}
    \begin{tabular}{|rllllp{11cm}|}\hline
        & & \multicolumn{3}{l}{} &
        \\
        & \multicolumn{4}{l}{
            \textbf{CatCode}
        }
        & \textit{\textbf{CatCode directory}}
        \\
        & > & \multicolumn{3}{l}{
            \textbf{Classes} 
        }
        & \textit{\textbf{Contains all of the super-classes for CatCode}}
        \\
        & & -- & \multicolumn{2}{l}{
            boolean\_operator.gd 
        }
        & \textit{Super class for all boolean operators}
        \\
        & & -- & \multicolumn{2}{l}{
            conditional\_block.gd 
        }
        & \textit{Super class for selective blocks}
        \\
        & & -- & \multicolumn{2}{l}{
            function\_block.gd 
        }
        & \textit{Super class for generalised functions}
        \\
        & & -- & \multicolumn{2}{l}{
            instruction.gd 
        }
        & \textit{Class for instruction lines}
        \\
        & & -- & \multicolumn{2}{l}{
            iteration\_block.gd 
        }
        & \textit{Super class for iterative blocks}
        \\
        & & -- & \multicolumn{2}{l}{
            logic\_block.gd 
        }
        & \textit{Super class for catCode logical blocks}
        \\
        & & -- & \multicolumn{2}{l}{
            parameter\_component.gd 
        }
        & \textit{Super class for parameter block components}
        \\
        & > & \multicolumn{3}{l}{
            \textbf{graphicalBlocks} 
        }
        & \textit{\textbf{Contains the UI blocks for Graphical CatCode}}
        \\
        & & -- & \multicolumn{2}{l}{
            boolean\_operand.gd/tscn
        }
        & \textit{Basic component for representing operands graphically}
        \\
        & & -- & \multicolumn{2}{l}{
            compoundOperand.gd
        }
        & \textit{Used for operand blocks, such as OP\_1\_operandContainer}
        \\
        & & -- & \multicolumn{2}{l}{
            graphical\_block.gd/tscn
        }
        & \textit{Generic class/node for graphical blocks}
        \\
        & & -- & \multicolumn{2}{l}{
            GB\_1\_operand.gd/tscn
        }
        & \textit{Generic class/node for graphical blocks with one parameter}
        \\
        & & -- & \multicolumn{2}{l}{
            GB\_\textbf{*}.tscn
        }
        & \textit{Node family for graphical block implementations}
        \\
        & & -- & \multicolumn{2}{l}{
            OP\_\textbf{*}.tscn
        }
        & \textit{Node family for graphical block implementation of operands}
        \\
        & > & \multicolumn{3}{l}{
            \textbf{graphicalComponents} 
        }
        & \textit{\textbf{Contains building blocks for graphical blocks}}
        \\
        & & -- & \multicolumn{2}{l}{
            basic\_indent\_node.tscn
        }
        & \textit{Practically empty node for indenting Graphical CatCode}
        \\
        & & -- & \multicolumn{2}{l}{
            block\_auto\_identifier.gd/tscn
        }
        & \textit{A label that will automatically update to a block's name}
        \\
        & & -- & \multicolumn{2}{l}{
            \textbf{*}\_param.gd/tscn
        }
        & \textit{Node family for GUI representation of parameters}
        \\
        & & -- & \multicolumn{2}{l}{
            operand\_container.gd/tscn
        }
        & \textit{Node for containing operand nodes}
        \\
        & > & \multicolumn{3}{l}{
            \textbf{logicBlocks} 
        }
        & \textit{\textbf{Contains the logical, back-end nodes for CatCode}}
        \\
        & > & > & \multicolumn{2}{l}{
            \textbf{nodes} 
        }
        & \textit{\textbf{Contains the nodes/scenes for the logic blocks}}
        \\
        & & & -- & \multicolumn{1}{l}{
            \textbf{*}.tscn
        }
        & \textit{View \autoref{tab:catCodeSpec/logic} and \autoref{tab:catCodeSpec/op} for file associations}
        \\
        & > & > & \multicolumn{2}{l}{
            \textbf{scripts} 
        }
        & \textit{\textbf{Contains scripts for scripts for logic blocks}}
        \\
        & & & -- & \multicolumn{1}{l}{
            \textbf{*}.gd
        }
        & \textit{View \autoref{tab:catCodeSpec/logic} and \autoref{tab:catCodeSpec/op} for file associations}
        \\
        & > & \multicolumn{3}{l}{
            \textbf{preconfigCCInstances} 
        }
        & \textit{\textbf{Contains codeManagers with blocks already added}}
        \\
        & & \multicolumn{3}{l}{...} &
        \\
        & & \multicolumn{3}{l}{} &
        \\
    \hline
    \end{tabular}
    \label{fig:catCodeSpec/files}
\end{fullwidthfigure}

\begin{fullwidth}
\centering
\begin{tabular}{|rllllp{11cm}|}\hline
    & & \multicolumn{3}{l}{} &
    \\
    & & \multicolumn{3}{l}{...} &
    \\
    & > & \multicolumn{3}{l}{
        \textbf{ui} 
    }
    & \textit{\textbf{Folder for all the GUI elements}}
    \\
    & & -- & \multicolumn{2}{l}{
        basic\_graphical\_editor.gd/tscn
    }
    & \textit{Basic graphical cat code editor node}
    \\
    & & -- & \multicolumn{2}{l}{
        basic\_output.gd/tscn
    }
    & \textit{Basic output node}
    \\
    & & -- & \multicolumn{2}{l}{
        basic\_textual\_editor.gd/tscn
    }
    & \textit{Basic editor for text-based CatCode}
    \\
    & -- & \multicolumn{3}{l}{
        cat\_code\_manager.gd/tscn
    }
    & \textit{Manages catCode and controls the execution of it}
    \\
    & -- & \multicolumn{3}{l}{
        cat\_thread.gd/tscn
    }
    & \textit{The seperate runner for CatCode}
    \\
    & & \multicolumn{3}{l}{} &
    \\
\hline
\end{tabular}
\end{fullwidth}

\section{Logic Blocks}

\begin{center}
  \begin{longtable}{lllr}
    \caption{CatCode Specification: Logic blocks} \label{tab:catCodeSpec/logic} \\
    \toprule
    \textbf{Name}
    & \texttt{Reference}
    & \texttt{Category}
    & \textit{Parameters} \\
    \multicolumn{4}{l}{
        Description
    }
    \\\midrule
    \endfirsthead
    \toprule
    \textbf{Name}
    & \texttt{Reference}
    & \texttt{Category}
    & \textit{Parameters} \\
    \multicolumn{4}{l}{
        Description
    }
    \\\midrule
    \endhead
    \textbf{If}
    & \texttt{if}
    & \texttt{conditional}
    & \textit{Condition (boolean operator)}
    \\
    \multicolumn{4}{p{0.98\linewidth}}{
        SELECTIVE BLOCK.
        Contains a boolean operator (view operator section).
        The section indented beneath this block will be executed if the operator is true.
        If the statement is not true, it will skip the section.
    }
    \\ \multicolumn{4}{l}{Node: \texttt{if\_conditional.tcsn}}
    \\ \multicolumn{4}{l}{Script: \texttt{if\_conditional.gd}}
    \\ \multicolumn{4}{l}{UI: \texttt{GB\_1\_Operand.tcsn}}
    \\ \\
    
    \textbf{Else if}
    & \texttt{elif}
    & \texttt{conditional}
    & \textit{Condition (boolean operator)}
    \\
    \multicolumn{4}{p{0.98\linewidth}}{
        SELECTIVE BLOCK.
        An if statement that will only be considered if none of the previous selective blocks are true.
    }
    \\ \multicolumn{4}{l}{Node: \texttt{else\_if\_conditional.tcsn}}
    \\ \multicolumn{4}{l}{Script: \texttt{if\_conditional.gd}}
    \\ \multicolumn{4}{l}{UI: \texttt{GB\_1\_Operand.tcsn}}
    \\ \\
    
    \textbf{Else}
    & \texttt{else}
    & \texttt{conditional}
    & \textit{--}
    \\
    \multicolumn{4}{p{0.98\linewidth}}{
        SELECTIVE BLOCK.
        The section indented after this block will only be executed if none of the previous selective blocks are true.
    }
    \\ \multicolumn{4}{l}{Node: \texttt{else\_conditional.tcsn}}
    \\ \multicolumn{4}{l}{Script: \texttt{if\_conditional.gd}}
    \\ \multicolumn{4}{l}{UI: \texttt{graphical\_block.tcsn}}
    \\ \\
    
    \textbf{Jump}
    & \texttt{jump}
    & \texttt{action}
    & \textit{Strength (-1.0 < float < 1.0)}
    \\
    \multicolumn{4}{p{0.98\linewidth}}{
        Causes catbot to jump.
        "Strength" controls the height of the job.
        Will result in an error if used when CatBot isn't on the ground.
    }
    \\ \multicolumn{4}{l}{Node: \texttt{jump\_function.tcsn}}
    \\ \multicolumn{4}{l}{Script: \texttt{jump\_function.gd}}
    \\ \multicolumn{4}{l}{UI: \texttt{GB\_1\_StrengthParam.tcsn}}
    \\ \\

    \textbf{Walk}
    & \texttt{moveX}
    & \texttt{action}
    & \textit{Strength (-1.0 < float < 1.0)}
    \\
    \multicolumn{4}{p{0.98\linewidth}}{
        Causes catbot to walk in a direction.
        Strength =  100\% Full speed to the right,
        Strength = -100\% Full speed to the left.
    }
    \\ \multicolumn{4}{l}{Node: \texttt{move\_x\_function.tcsn}}
    \\ \multicolumn{4}{l}{Script: \texttt{move\_x\_function.gd}}
    \\ \multicolumn{4}{l}{UI: \texttt{GB\_1\_StrengthParam.tcsn}}
    \\ \\

    \textbf{Walk}
    & \texttt{moveX}
    & \texttt{action}
    & \textit{Strength (-1.0 < float < 1.0)}
    \\
    \multicolumn{4}{p{0.98\linewidth}}{
        Causes catbot to walk in a direction.
        Strength =  100\% Full speed to the right,
        Strength = -100\% Full speed to the left.
    }
    \\ \multicolumn{4}{l}{Node: \texttt{move\_x\_function.tcsn}}
    \\ \multicolumn{4}{l}{Script: \texttt{move\_x\_function.gd}}
    \\ \multicolumn{4}{l}{UI: \texttt{GB\_1\_StrengthParam.tcsn}}
    \\ \\

    \textbf{Pass}
    & \texttt{pass}
    & \texttt{debug}
    & \textit{--}
    \\
    \multicolumn{4}{p{0.98\linewidth}}{
        Does literally nothing.
    }
    \\
    \multicolumn{4}{p{0.98\linewidth}}{
        \textit{Added for backend and testing reasons, left player facing as a joke.}
    }
    \\ \multicolumn{4}{l}{Node: \texttt{pass\_function.tcsn}}
    \\ \multicolumn{4}{l}{Script: \texttt{pass\_function.gd}}
    \\ \multicolumn{4}{l}{UI: \texttt{graphical\_block.tcsn}}
    \\ \\

    \textbf{Print}
    & \texttt{print}
    & \texttt{output}
    & \textit{Text (String)}
    \\
    \multicolumn{4}{p{0.98\linewidth}}{
        Outputs the contents of the parameter to CatBot's console.
    }
    \\ \multicolumn{4}{l}{Node: \texttt{print\_line\_function.tcsn}}
    \\ \multicolumn{4}{l}{Script: \texttt{print\_line\_function.gd}}
    \\ \multicolumn{4}{l}{UI: \texttt{GB\_1\_textBox.tcsn}}
    \\ \\

    \textbf{Print2}
    & \texttt{print2}
    & \texttt{output}
    & \textit{Text1 (String), Text2 (String)}
    \\
    \multicolumn{4}{p{0.98\linewidth}}{
        Outputs the contents of the parameters to CatBot's console.
    }
    \\
    \multicolumn{4}{p{0.98\linewidth}}{
        \textit{Added to test having multiple parameters in one function.}
    }
    \\ \multicolumn{4}{l}{Node: \texttt{print\_2\_function.tcsn}}
    \\ \multicolumn{4}{l}{Script: \texttt{print\_2\_function.gd}}
    \\ \multicolumn{4}{l}{UI: \texttt{GB\_2\_textBox.tcsn}}
    \\ \\

    \textbf{Repeat}
    & \texttt{repeat}
    & \texttt{loops}
    & \textit{Iterations (Integer > 0)}
    \\
    \multicolumn{4}{p{0.98\linewidth}}{
        ITERATIVE BLOCK
        The code section indented beneath this will be repeated the specified number of times.
        One iteration will occur per frame.
        If this is getting used, you may need to manually save the code for CatBot's behaviour to update.
    }
    \\ \multicolumn{4}{l}{Node: \texttt{repeat\_iterational.tcsn}}
    \\ \multicolumn{4}{l}{Script: \texttt{repeat\_iterational.gd}}
    \\ \multicolumn{4}{l}{UI: \texttt{GB\_1\_openPossitiveIntBox.tcsn}}
    \\ \\

    \textbf{Repeat Forever}
    & \texttt{repeatForever}
    & \texttt{loops}
    & \textit{--}
    \\
    \multicolumn{4}{p{0.98\linewidth}}{
        ITERATIVE BLOCK
        The code section indented beneath this will be repeated forever.
        One iteration will occur per frame.
        If this is getting used, you will need to manually save the code for CatBot's behaviour to update.
    }
    \\ \multicolumn{4}{l}{Node: \texttt{repeat\_forever\_iterational.tcsn}}
    \\ \multicolumn{4}{l}{Script: \texttt{repeat\_forever\_iterational.gd}}
    \\ \multicolumn{4}{l}{UI: \texttt{GB\_1\_fakeParam.tcsn}}
    \\ \\

    \textbf{Stripe}
    & \texttt{swipe}
    & \texttt{action}
    & \textit{--}
    \\
    \multicolumn{4}{p{0.98\linewidth}}{
        Causes CatBot to swipe.
        CatBot can swipe twice in midair.
        If this number is exceeded, an error will occur.
    }
    \\ \multicolumn{4}{l}{Node: \texttt{strike\_function.tcsn}}
    \\ \multicolumn{4}{l}{Script: \texttt{strike\_function.gd}}
    \\ \multicolumn{4}{l}{UI: \texttt{GB\_1\_fakeParam.tcsn}}
    \\ \\
    \bottomrule
  \end{longtable}
\end{center}

\section{Boolean Operators}

\begin{center}
  \begin{longtable}{lllr}
    \caption{CatCode Specification: Boolean operators} \label{tab:catCodeSpec/op} \\
    \toprule
    \textbf{Name}
    & \texttt{Reference}
    & \texttt{Category}
    & \textit{Parameters} \\
    \multicolumn{4}{l}{
        Description
    }
    \\\midrule
    \endfirsthead
    \toprule
    \textbf{Name}
    & \texttt{Reference}
    & \texttt{Category}
    & \textit{Parameters} \\
    \multicolumn{4}{l}{
        Description
    }
    \\\midrule
    \endhead
    
    \textbf{And}
    & \texttt{and}
    & \texttt{logic}
    & \textit{Op1, Op2 (Boolean Operators)}
    \\
    \multicolumn{4}{p{0.98\linewidth}}{
        Contains two operators.
        Returns true if both operators are true.
        This means that code indented beneath an if statement executes if both statements are true.
    }
    \\
    \multicolumn{4}{p{0.98\linewidth}}{
        \textit{Not compatible with graphical CatCode.}
    }
    \\ \multicolumn{4}{l}{Node: \texttt{and\_operand.tcsn}}
    \\ \multicolumn{4}{l}{Script: \texttt{and\_operand.gd}}
    \\ \multicolumn{4}{l}{UI: \texttt{OP\_2\_operandContainers.tcsn}}
    \\ \\

    \textbf{<True/False>}
    & \texttt{<true/false>}
    & \texttt{value}
    & \textit{--}
    \\
    \multicolumn{4}{p{0.98\linewidth}}{
        Represents a boolean value of <true/false>.
    }
    \\
    \multicolumn{4}{p{0.98\linewidth}}{
        \textit{This block is used for True and False values. The exact value is set in the Godot editor before runtime and is two separate blocks as far as the player is concerned.}
    }
    \\ \multicolumn{4}{l}{Node: \texttt{boolean\_value.tcsn}}
    \\ \multicolumn{4}{l}{Script: \texttt{boolean\_value.gd}}
    \\ \multicolumn{4}{l}{UI: \texttt{basic\_operand.tcsn}}
    \\ \\

    \textbf{Can Swipe}
    & \texttt{canSwipe}
    & \texttt{physics logic}
    & \textit{--}
    \\
    \multicolumn{4}{p{0.98\linewidth}}{
        Will return true if CatBot can swipe.
    }
    \\ \multicolumn{4}{l}{Node: \texttt{can\_strike\_operand.tcsn}}
    \\ \multicolumn{4}{l}{Script: \texttt{can\_strike\_operand.gd}}
    \\ \multicolumn{4}{l}{UI: \texttt{basic\_operand.tcsn}}
    \\ \\

    \textbf{Not}
    & \texttt{not}
    & \texttt{logic}
    & \textit{Operand (Boolean Operator)}
    \\
    \multicolumn{4}{p{0.98\linewidth}}{
        Contains one operator.
        Returns true if the operator is false.
        True becomes false, and false becomes true, for example.
    }
    \\
    \multicolumn{4}{p{0.98\linewidth}}{
        \textit{Unknown compatibility with graphical CatCode.}
    }
    \\ \multicolumn{4}{l}{Node: \texttt{not\_operand.tcsn}}
    \\ \multicolumn{4}{l}{Script: \texttt{not\_operand.gd}}
    \\ \multicolumn{4}{l}{UI: \texttt{OP\_1\_operandContainer.tcsn}}
    \\ \\
    
    \textbf{Or}
    & \texttt{or}
    & \texttt{logic}
    & \textit{Op1, Op2 (Boolean Operators)}
    \\
    \multicolumn{4}{p{0.98\linewidth}}{
        Contains two operators.
        Returns true if either operator is true.
        This means that code indented beneath an if statement executes if one statement is true.
    }
    \\
    \multicolumn{4}{p{0.98\linewidth}}{
        \textit{Not compatible with graphical CatCode.}
    }
    \\ \multicolumn{4}{l}{Node: \texttt{or\_operand.tcsn}}
    \\ \multicolumn{4}{l}{Script: \texttt{or\_operand.gd}}
    \\ \multicolumn{4}{l}{UI: \texttt{OP\_2\_operandContainers.tcsn}}
    \\ \\
    
    \bottomrule
  \end{longtable}
\end{center}

\break
\section{Help Messages}

These help messages are brought up by pressing the associated button the graphical CatCode editor. \textquote{Instruction Set}, which brings up an automatically generated list of the current instruction set of blocks and their descriptions, is not shown here.

\begin{figure}[h]
    \centering
    \begin{tcolorbox}[boxrule=0.25mm]
        CatCode+ is the graphical code editor, that created CatCode which controls the CatBot. \\
        \phantom{-} \\
        \textbf{What is catCode?} \\
        CatCode consists of a sequence of blocks. Each of these are an instruction that the catBot will follow in order. For example, the "jump" block will make catBot execute the jump function. \\
        \phantom{-} \\
        Some blocks have values associated with them called parameters. When these functions are executed, the value specified by the parameter (or sometimes parameters) will affect how the catBot will react. For example, the strength parameter on "jump" will control how high the catBot will jump. \\
        \phantom{-} \\
        Experiment with these! Unlike CatCode-, it is hard to give catBot functions it can't handle using catCode+. \\
        \phantom{-} \\
        To view a list of instructions, their parameters and further explanations of the blocks click "Instruction set". \\
        \phantom{-} \\
        \textbf{How do I edit CatCode?} \\
        To add a new block, click on the "+ add line" button. To remove a block, right click it. Click and drag blocks to reorder them. \\
        \phantom{-} \\
        \textbf{How do I add indent blocks?} \\
        To add indent blocks, use the "<" and ">" buttons on it's line. \\
        \phantom{-} \\
        \textbf{Why do I indent blocks?} \\
        Certain blocks, such as "if" will modify the behaviour of the blocks associated with it. Indent blocks are used to tell catBot what is beneath what. \\
        \phantom{-} \\
        When inserting a new line, the block will be inserted at the correct indent. \\
        \phantom{-} \\
        \textit{Click outside window to close}
    \end{tcolorbox}
    \caption[Help pop-up text]{"Help" pop-up text, taken from survey build}
    \label{fig:catCodeSpec/help}
\end{figure}

\begin{figure}[h]
    \centering
    \begin{tcolorbox}[boxrule=0.25mm]
        At present controls are hardwired.\\
        If you need to remap controls, return to the coverpage and follow the unexported build instructions. These instructions will guide you through the process of installing godot, downloading the project and editing the controls. This will \textit{not} update catCode labelling.\\
        \phantom{-} \\
        \textbf{Controls}\\
        Z-Button: \textit{Z}\\
        X-Button: \textit{X}\\
        \phantom{-} \\
        Up: \textit{Up arrow key}\\
        Down: \textit{Down arrow key}\\
        left: \textit{Left arrow key}\\
        Right: \textit{Right arrow key}\\
        \phantom{-} \\
        Delete block: \textit{Right click}\\
        \phantom{-} \\
        It is up to you what each of the above keys do, however when relying on the autocontrollers, the following keybinds are used:\\
        \phantom{-} \\
        \textbf{Auto keybindings}\\
        Z-Button: \textit{Jump}\\
        X-Button: \textit{Strike/Swipe}\\
        Arrow keys: \textit{Movement}\\
        \phantom{-} \\
        \textit{Click outside window to close}
    \end{tcolorbox}
    \caption[Controls pop-up text]{"Controls" pop-up text, taken from survey build}
    \label{fig:catCodeSpec/control}
\end{figure}